% Generated draft: requires mathematical and source review.
\documentclass[11pt]{article}
\usepackage[margin=25mm]{geometry}
\usepackage{fontspec}
\setmainfont{Latin Modern Roman}
\usepackage{amsmath,amssymb,mathtools,unicode-math}
\setmathfont{Latin Modern Math}
\usepackage{enumitem,needspace,xurl,hyperref}
\setlength{\parindent}{0pt}
\setlength{\parskip}{5pt}
\setlength{\emergencystretch}{3em}
\newcommand{\N}{\mathbb{N}}
\newcommand{\Z}{\mathbb{Z}}
\newcommand{\Q}{\mathbb{Q}}
\newcommand{\R}{\mathbb{R}}
\newcommand{\C}{\mathbb{C}}
\newcommand{\F}{\mathbb{F}}
\newcommand{\A}{\mathbb{A}}
\newcommand{\PP}{\mathbb{P}}
\newcommand{\E}{\mathbb{E}}
\newcommand{\eps}{\varepsilon}
\newcommand{\dd}{\,\mathrm{d}}
\begin{document}
\section*{UnsolvedMath ID: 6700014}
\textbf{Source problem code:} catalogue entry
\subsection{Definitions and mathematical statement}
For a smooth manifold \(X\), let
\[
\begin{split}
\mathcal R_0(X)&=\{g:g\text{ is a smooth Riemannian metric on }X,\ \operatorname{Ric}_g=0\},\\
\mathcal R_+(X)&=\{g:g\text{ is a smooth Riemannian metric on }X,\ \operatorname{Scal}_g(x)>0\text{ for all }x\in X\}.
\end{split}
\]
Gromov asks how common Ricci-flat metrics are on compact simply connected manifolds \(X\) for which \(\mathcal R_+(X)\ne\varnothing\). In particular, the existence aspect concerns those \(X\) satisfying
\[
\pi_1(X)=0,\qquad \mathcal R_+(X)\ne\varnothing,\qquad \mathcal R_0(X)\ne\varnothing.
\]
The two nonemptiness conditions concern potentially different metrics on the same smooth manifold. A Ricci-flat metric has zero scalar curvature, so it cannot itself belong to \(\mathcal R_+(X)\).

The phrase ``how common'' has no supplied counting rule, topology, measure, or moduli-space convention. It may ask for the prevalence of such manifolds, or for the size of \(\mathcal R_0(X)\) on a fixed admissible \(X\); these are different targets. No dimension or boundary convention is given. Thus the existence condition can be stated precisely, but it does not by itself replace the broader prevalence question.
\subsection{Short English statement}
On compact simply connected manifolds that support positive scalar curvature metrics, investigate the existence and prevalence of Ricci-flat metrics. The two curvature conditions apply to different metrics, and the meaning of prevalence remains unspecified.
\subsection{Sources}
\begin{itemize}
\item Local UnsolvedMath registry, permanent ID 6700014, source code catalogue entry; Misha Gromov, 2017; update 2026-08-06.
\item Gromov, \emph{101 Questions, Problems and Conjectures around Scalar Curvature} (2017), printed tag [?14], occurrence 14, PDF page 15: \url{https://www.ihes.fr/~gromov/wp-content/uploads/2018/08/101-problemsOct1-2017.pdf}. Registry provenance reports access on 2026-07-29 and normalized PDF extraction from an incomplete public author-hosted PDF. Present status and reuse terms need review.
\item The supplied research summary discusses Calabi--Yau manifolds, rigidity, and a putative K3 example without references. It contains conflicting assertions about simple connectedness and conflates simultaneous curvature conditions on a metric with coexistence of different metrics. Its examples and status assertions require independent checking and are not hypotheses of this draft.
\end{itemize}
\end{document}
